\subsection{$A_1$ 权与一个外推定理}
\label{sec:ch7-a1-extrapolation}

本节先利用 Hardy--Littlewood 极大函数给出 $A_1$ 的构造性刻画. 结合命题~\ref{prop:ch7-ap-properties}, 这使我们能对所有 $p$ 构造 $A_p$ 权. 随后用这一构造证明一个强有力的外推定理.

\begin{Theorem}
\label{thm:ch7-a1-characterization}
有下列结论:
\begin{enumerate}[label=(\arabic*)]
\item 设 $f\in L^1_{\mathrm{loc}}(\mathbb R^n)$ 且 $Mf(x)<\infty$ 几乎处处成立. 若 $0<\delta<1$, 则 $w(x)=Mf(x)^\delta$ 是 $A_1$ 权, 其 $A_1$ 常数只依赖于 $\delta$.
\item 反之, 若 $w\in A_1$, 则存在 $f\in L^1_{\mathrm{loc}}(\mathbb R^n)$, $\delta<1$ 以及满足 $K,K^{-1}\in L^\infty$ 的 $K$, 使 $w(x)=K(x)Mf(x)^\delta$.
\end{enumerate}
\end{Theorem}

\begin{Proof}
为证第一条, 只需证明存在常数 $C$, 使对任意 $f$, 任意立方体 $Q$ 和几乎每个 $x\in Q$,
\[
 \frac1{|Q|}\int_Q(Mf)^\delta\le CMf(x)^\delta.
\]
固定 $Q$, 写 $f=f_1+f_2$, 其中 $f_1=f\chi_{2Q}$. 则 $Mf\le Mf_1+Mf_2$, 因而
\[
 Mf(x)^\delta\le Mf_1(x)^\delta+Mf_2(x)^\delta.
\]
由 $M$ 的弱型 $(1,1)$ 性质和 Kolmogorov 不等式,
\[
 \frac1{|Q|}\int_QMf_1(y)^\delta\,dy
 \le\frac{C_n}{|Q|^\delta}\left(\int_{2Q}|f|\right)^\delta
 \le C_nMf(x)^\delta.
\]

\hypertarget{chapter-07-page-153}{}
为估计 $Mf_2$, 注意若 $y\in Q$, $R$ 是包含 $y$ 且 $\int_R|f_2|>0$ 的立方体, 则其边长满足 $\ell(R)>\ell(Q)/2$. 因而存在只依赖于 $n$ 的 $c_n$, 使对 $x\in Q$ 有 $x\in c_nR$. 所以
\[
 \frac1{|R|}\int_R|f_2|
 \le\frac{c_n^n}{|c_nR|}\int_{c_nR}|f_2|
 \le c_n^nMf(x).
\]
故对任意 $y\in Q$, $Mf_2(y)\le c_n^nMf(x)$, 从而
\[
 \frac1{|Q|}\int_QMf_2(y)^\delta\,dy\le c_n^{n\delta}Mf(x)^\delta.
\]

若 $w\in A_1$, 则由反向 Hölder 不等式, 存在 $\varepsilon>0$ 使
\[
 \left(\frac1{|Q|}\int_Qw^{1+\varepsilon}\right)^{1/(1+\varepsilon)}
 \le\frac C{|Q|}\int_Qw.
\]
结合 $A_1$ 条件可得
\[
 M(w^{1+\varepsilon})(x)^{1/(1+\varepsilon)}
 \le Cw(x)
 \quad\text{几乎处处}.
\]
Lebesgue 微分定理给出常数为 $1$ 的反向不等式. 令 $f=w^{1+\varepsilon}$, $\delta=1/(1+\varepsilon)$, 得
\[
 w(x)\le Mf(x)^\delta\le Cw(x).
\]
取 $K(x)=w(x)/Mf(x)^\delta$, 即得第二条.
\end{Proof}

第一条中的 $f$ 还可替换为满足 $M\mu(x)<\infty$ 几乎处处的有限 Borel 测度 $\mu$. 特别地, 若 $\delta_0$ 是原点处的 Dirac 测度, 则 $M\delta_0(x)=C|x|^{-n}$. 因而当 $-n<\alpha\le0$ 时 $|x|^\alpha\in A_1$; 由命题~\ref{prop:ch7-ap-properties}, 当 $-n<\alpha<n(p-1)$ 时 $|x|^\alpha\in A_p$. 该范围是最佳的, 因为在范围外 $|x|^\alpha$ 与 $|x|^{\alpha(1-p')}$ 不会同时局部可积.

\begin{Theorem}
\label{thm:ch7-rubio-de-francia-extrapolation}
固定 $1<r<\infty$. 假设对任意 $w\in A_r$, 算子 $T$ 在 $L^r(w)$ 上有界, 且算子范数只依赖于 $w$ 的 $A_r$ 常数. 则对每个 $1<p<\infty$ 和每个 $w\in A_p$, $T$ 在 $L^p(w)$ 上有界.
\end{Theorem}

\begin{Proof}
先证明若 $1<q<r$ 且 $w\in A_1$, 则 $T$ 在 $L^q(w)$ 上有界. 由定理~\ref{thm:ch7-a1-characterization}, $(Mf)^{(r-q)/(r-1)}\in A_1$. 再由命题~\ref{prop:ch7-ap-properties}, $w(Mf)^{q-r}$ 是 $A_r$ 权. 因而由 Hölder 不等式, 假设以及定理~\ref{thm:ch7-maximal-weighted-strong},
\[
\begin{aligned}
 \int_{\mathbb R^n}|Tf|^qw
 &\le\left(\int_{\mathbb R^n}|Tf|^rw(Mf)^{q-r}\right)^{q/r}
 \left(\int_{\mathbb R^n}(Mf)^qw\right)^{(r-q)/r}\\
 &\le C\left(\int_{\mathbb R^n}|f|^rw(Mf)^{q-r}\right)^{q/r}
 \left(\int_{\mathbb R^n}|f|^qw\right)^{(r-q)/r}\\
 &\le C\int_{\mathbb R^n}|f|^qw.
\end{aligned}
\]
最后一步使用 $|f|\le Mf$ 以及 $q-r<0$.

\hypertarget{chapter-07-page-154}{}
现在证明: 给定 $1<p<\infty$ 与 $1<q<\min(p,r)$, 若 $w\in A_{p/q}$, 则 $T$ 在 $L^p(w)$ 上有界. 这已经足够, 因为给定 $w\in A_p$, 推论~\ref{cor:ch7-ap-self-improvement} 给出某个 $q>1$ 使 $w\in A_{p/q}$.

固定 $w\in A_{p/q}$. 由对偶性, 存在范数为 $1$ 的 $u\in L^{(p/q)'}(w)$, 使
\[
 \left(\int_{\mathbb R^n}|Tf|^pw\right)^{q/p}
 =\int_{\mathbb R^n}|Tf|^qwu.
\]
对任意 $s>1$, 有 $wu\le M((wu)^s)^{1/s}$, 且 $M((wu)^s)^{1/s}\in A_1$. 由证明的第一部分,
\[
\begin{aligned}
 \int_{\mathbb R^n}|Tf|^qwu
 &\le\int_{\mathbb R^n}|Tf|^qM((wu)^s)^{1/s}\\
 &\le C\int_{\mathbb R^n}|f|^qM((wu)^s)^{1/s}\\
 &\le C\left(\int_{\mathbb R^n}|f|^pw\right)^{q/p}
 \left(\int_{\mathbb R^n}M((wu)^s)^{(p/q)'/s}
 w^{1-(p/q)'}\right)^{1/(p/q)'}.
\end{aligned}
\]
由 $w\in A_{p/q}$ 和命题~\ref{prop:ch7-ap-properties}, $w^{1-(p/q)'}\in A_{(p/q)'}$. 取 $s$ 充分接近 $1$, 使 $w^{1-(p/q)'}\in A_{(p/q)'/s}$. 于是定理~\ref{thm:ch7-maximal-weighted-strong} 表明第二个积分由
\[
 C\int_{\mathbb R^n}(wu)^{(p/q)'}w^{1-(p/q)'}=C
\]
控制, 证明完成.
\end{Proof}

同样的论证给出如下变形: 给定 $s\ge1$, 若对所有 $w\in A_{r/s}$, $T$ 在 $L^r(w)$ 上有界, 则对 $p>s$ 和所有 $w\in A_{p/s}$, $T$ 在 $L^p(w)$ 上有界.
